\documentclass[12pt]{ctexart}
\usepackage{graphicx}
\usepackage{float}
\graphicspath{{../figures/}}

\title{单问最小示例：线性趋势基线模型}
\author{}
\date{}

\begin{document}
\maketitle

\section{问题分析}

问题一要求根据某站点过去 5 天需求量预测第 6 天需求量。该最小示例只演示单问求解时的文件形态，因此采用可手算、可解释的线性趋势基线模型。即使是最小示例，也要把建模、求解、验证和结论写回 \verb|paper/main.tex|。

\section{模型建立与求解}

设第 \(t\) 天需求量为 \(y_t\)。由第 \(1\) 天至第 \(5\) 天需求量从
\(80\) 件增加到 \(99\) 件，可得平均日增长量
\[
k=\frac{99-80}{5-1}=4.75.
\]
因此第 \(6\) 天预测需求量为
\[
\hat y_6=y_5+k=103.75\text{ 件}.
\]
结果保存于 \verb|tables/tab_q1_result.csv|。

\section{结果与验证}

图 \verb|fig_q1_model_flow.svg| 给出最终模型流程，流程图源文件保存在
\verb|modeling/q1_model_flow.mmd|，图
\verb|fig_q1_result.svg| 展示历史需求与预测结果。由于样本只有 5 天，
该结果应作为基线预测，而不能包装成高精度模型。正式竞赛中应补充更多
历史数据，并与移动平均、季节性基线或滚动验证结果比较。

\section{结论}

在当前最小数据下，第 \(6\) 天需求量的线性趋势预测值为
\(103.75\) 件。该单问示例的论文成品只使用 \verb|paper/main.tex|，
不再拆分为 \verb|paper/sections/q1.tex|。

\end{document}
